% =====================================================================
%  RESUMEN EJECUTIVO (castellano / valenciano / inglés)
%  NOTA: la versión en valenciano debería ser revisada por un hablante nativo.
% =====================================================================
\chapter*{Resumen ejecutivo}
\addcontentsline{toc}{chapter}{Resumen ejecutivo}
\markboth{Resumen ejecutivo}{Resumen ejecutivo}

\section*{Resumen}

En el marco del Pacto Verde Europeo y de la creciente presión normativa y social, la
sostenibilidad ha dejado de ser un tema marginal en la comunicación corporativa: lo que está en
juego no es solo sus prácticas ambientales y sociales, sino también la manera en que las
comunican. Instrumentos legislativos como la Directiva sobre informes de sostenibilidad
corporativa (CSRD) y los estándares ESRS, la Taxonomía Europea y el Reglamento SFDR, que
sustituyen y amplían el régimen previo de la NFRD, han redefinido las expectativas sobre la
transparencia y la rendición de cuentas. En este contexto, el presente trabajo analiza cómo las
grandes empresas europeas comunican sus estrategias de sostenibilidad a través de sus informes
corporativos, centrándose en la información de sostenibilidad integrada en el informe de
gestión, que es donde la normativa europea obliga a incluirla y que constituye una pieza clave
del relato empresarial.

Para llevar a cabo este análisis se aplican técnicas de inteligencia artificial, concretamente
de procesamiento del lenguaje natural (PLN), que permiten examinar de manera objetiva,
sistemática y escalable el contenido de los informes de una muestra de empresas del índice
STOXX Europe 600. El estudio combina modelización de temas, análisis de tono y especificidad
del lenguaje, diccionarios especializados (financieros y de los propios ESRS) y modelos de
lenguaje preentrenados, sin recurrir a calificaciones ESG de terceros. El objetivo es identificar
los principales discursos, enfoques estratégicos y patrones comunicativos empleados en la
presentación de la sostenibilidad, detectar posibles señales de \textit{greenwashing}, y examinar
su evolución ante el cambio de marco regulatorio, así como su relación con características
empresariales como el sector, el país y el tamaño, y su coherencia con la normativa de la Unión
Europea.

\medskip
\noindent\textbf{Palabras clave:} \textit{greenwashing}; CSRD; ESRS; procesamiento del lenguaje
natural; informes de sostenibilidad; STOXX Europe 600; análisis de contenido; ClimateBERT.

\section*{Resum}

En el marc del Pacte Verd Europeu i de la creixent pressió normativa i social, la sostenibilitat
ha deixat de ser un tema marginal en la comunicació corporativa: el que està en joc no és només
les seues pràctiques ambientals i socials, sinó també la manera en què les comuniquen.
Instruments legislatius com la Directiva sobre informes de sostenibilitat corporativa (CSRD) i
els estàndards ESRS, la Taxonomia Europea i el Reglament SFDR, que substitueixen i amplien el
règim previ de la NFRD, han redefinit les expectatives sobre la transparència i la rendició de
comptes. En aquest context, el present treball analitza com les grans empreses europees
comuniquen les seues estratègies de sostenibilitat a través dels seus informes corporatius,
centrant-se en la informació de sostenibilitat integrada en l'informe de gestió, que és on la
normativa europea obliga a incloure-la i que constitueix una peça clau del relat empresarial.

Per a dur a terme aquesta anàlisi s'apliquen tècniques d'intel·ligència artificial, concretament
de processament del llenguatge natural (PLN), que permeten examinar de manera objectiva,
sistemàtica i escalable el contingut dels informes d'una mostra d'empreses de l'índex STOXX
Europe 600. L'estudi combina modelització de temes, anàlisi de to i especificitat del llenguatge,
diccionaris especialitzats (financers i dels propis ESRS) i models de llenguatge preentrenats,
sense recórrer a qualificacions ESG de tercers. L'objectiu és identificar els principals
discursos, enfocaments estratègics i patrons comunicatius emprats en la presentació de la
sostenibilitat, detectar possibles senyals de \textit{greenwashing}, i examinar la seua evolució
davant el canvi de marc regulador, així com la seua relació amb característiques empresarials com
el sector, el país i la grandària, i la seua coherència amb la normativa de la Unió Europea.

\medskip
\noindent\textbf{Paraules clau:} \textit{greenwashing}; CSRD; ESRS; processament del llenguatge
natural; informes de sostenibilitat; STOXX Europe 600; anàlisi de contingut; ClimateBERT.

\section*{Abstract}

Within the framework of the European Green Deal and growing regulatory and social pressure,
sustainability has ceased to be a marginal issue in corporate communication: what is at stake is
not only companies' environmental and social practices, but also the way in which they
communicate them. Legislative instruments such as the Corporate Sustainability Reporting
Directive (CSRD) and the ESRS standards, the European Taxonomy and the SFDR Regulation, which
replace and extend the previous NFRD regime, have redefined expectations regarding transparency
and accountability. Against this backdrop, this dissertation analyses how large European
companies communicate their sustainability strategies through their corporate reports, focusing
on the sustainability information embedded in the management report, where European regulation
requires it to be included and which constitutes a key piece of the corporate narrative.

To carry out this analysis, artificial intelligence techniques are applied, specifically natural
language processing (NLP), which make it possible to examine the content of the reports of a
sample of companies from the STOXX Europe 600 index in an objective, systematic and scalable
way. The study combines topic modelling, tone and language-specificity analysis, specialised
dictionaries (financial and ESRS-specific) and pre-trained language models, without resorting to
third-party ESG ratings. The aim is to identify the main discourses, strategic approaches and
communicative patterns used in the presentation of sustainability, to detect possible signals of
greenwashing, and to examine their evolution in the face of the regulatory shift, as well as
their relationship with corporate characteristics such as sector, country and size, and their
consistency with European Union regulation.

\medskip
\noindent\textbf{Keywords:} greenwashing; CSRD; ESRS; natural language processing; sustainability
reports; STOXX Europe 600; content analysis; ClimateBERT.

\clearpage
